\chapter{2017年上海卷（春）}

\section{填空题}

\begin{problem}
设集合 \(A=\{1,2,3\}\)，集合 \(B=\{3,4\}\)，则 \(A\cup B=\)\fillinblank{}.
\end{problem}

\begin{answer}
\(\{1,2,3,4\}\)
\end{answer}

\begin{solution}
并集包含属于集合 \(A\) 或属于集合 \(B\) 的所有元素，因此 \(A\cup B=\{1,2,3,4\}\).
\end{solution}

\begin{problem}
不等式 \(\lvert x-1\rvert<3\) 的解集为\fillinblank{}.
\end{problem}

\begin{answer}
\((-2,4)\)
\end{answer}

\begin{solution}
由绝对值不等式的定义，\(-3<x-1<3\)，所以 \(-2<x<4\). 因而解集为 \((-2,4)\).
\end{solution}

\begin{problem}
若复数 \(z\) 满足 \(2\overline{z}-1=3+6\i\)（\(\i\) 是虚数单位），则 \(z=\)\fillinblank{}.
\end{problem}

\begin{answer}
\(2-3\i\)
\end{answer}

\begin{solution}
设 \(z=u+v\i\)，其中 \(u\)，\(v\in\R\)，则 \(\overline z=u-v\i\). 代入原式得 \(2u-2v\i-1=3+6\i\)，比较实部和虚部可得 \(2u-1=3\)，\(-2v=6\)，即 \(u=2\)，\(v=-3\). 所以 \(z=2-3\i\).
\end{solution}

\begin{problem}
若 \(\cos\alpha=\dfrac13\)，则 \(\sin\left(\alpha-\dfrac{\pi}{2}\right)=\)\fillinblank{}.
\end{problem}

\begin{answer}
\(-\dfrac13\)
\end{answer}

\begin{solution}
利用诱导公式，\(\sin\left(\alpha-\dfrac{\pi}{2}\right)=-\cos\alpha\). 由 \(\cos\alpha=\dfrac13\)，得 \(\sin\left(\alpha-\dfrac{\pi}{2}\right)=-\dfrac13\).
\end{solution}

\begin{problem}
若关于 \(x\)，\(y\) 的方程组 \(\begin{cases}
x+2y=4,\\
3x+ay=6
\end{cases}\)
无解，则实数 \(a=\)\fillinblank{}.
\end{problem}

\begin{answer}
\(6\)
\end{answer}

\begin{solution}
方程组无解，说明两条直线平行且不重合，因此有 \(\dfrac31=\dfrac a2\)，而 \(\dfrac64\ne\dfrac31\). 由 \(\dfrac31=\dfrac a2\) 得 \(a=6\)，此时 \(\dfrac64=\dfrac32\ne3\)，确实不重合，所以 \(a=6\).
\end{solution}

\begin{problem}
若等差数列 \(\{a_n\}\) 的前 \(5\) 项的和为 \(25\)，则 \(a_1+a_5=\)\fillinblank{}.
\end{problem}

\begin{answer}
\(10\)
\end{answer}

\begin{solution}
等差数列前 \(5\) 项的和满足 \(S_5=\dfrac{5(a_1+a_5)}2\). 由 \(S_5=25\)，得 \(a_1+a_5=10\).
\end{solution}

\begin{problem}
若 \(P\)，\(Q\) 是圆 \(x^2+y^2-2x+4y+4=0\) 上的动点，则 \(\lvert PQ\rvert\) 的最大值为\fillinblank{}.
\end{problem}

\begin{answer}
\(2\)
\end{answer}

\begin{solution}
将圆的方程配方得 \((x-1)^2+(y+2)^2=1\)，所以圆心为 \((1,-2)\)，半径为 \(1\). 圆上两点距离的最大值是直径，故 \(\lvert PQ\rvert_{\max}=2\).
\end{solution}

\begin{problem}
已知数列 \(\{a_n\}\) 的通项公式为 \(a_n=3^n\)，则 \(\displaystyle\lim_{n\to\infty}\dfrac{a_1+a_2+\cdots+a_n}{a_n}=\)\fillinblank{}.
\end{problem}

\begin{answer}
\(\dfrac32\)
\end{answer}

\begin{solution}
由 \(a_n=3^n\)，
\[
\frac{a_1+a_2+\cdots+a_n}{a_n}=\frac{3+3^2+\cdots+3^n}{3^n}=\frac{3(3^n-1)}{2\cdot3^n}=\frac32\left(1-\frac1{3^n}\right)
\]
因为 \(3^{-n}\to0\)，所以所求极限为 \(\dfrac32\).
\end{solution}

\begin{problem}
若 \(\left(x+\dfrac{2}{x}\right)^n\) 的二项展开式的各项系数之和为 \(729\)，则该展开式中常数项的值为\fillinblank{}.
\end{problem}

\begin{answer}
\(160\)
\end{answer}

\begin{solution}
二项式各项系数之和等于令 \(x=1\) 时的值，因此 \(3^n=729=3^6\)，从而 \(n=6\). 通项为 \(\mathrm C_6^k x^{6-k}\left(\dfrac2x\right)^k=\mathrm C_6^k2^k x^{6-2k}\). 常数项要求 \(6-2k=0\)，故 \(k=3\)，常数项的值为 \(\mathrm C_6^3 2^3=20\times8=160\).
\end{solution}

\begin{problem}
设椭圆 \(\dfrac{x^2}{2}+y^2=1\) 的左，右焦点分别为 \(F_1\)，\(F_2\)，点 \(P\) 在该椭圆上，则使得 \(\triangle F_1F_2P\) 是等腰三角形的点 \(P\) 的个数是\fillinblank{}.
\end{problem}

\begin{answer}
\(6\)
\end{answer}

\begin{solution}
椭圆的长半轴平方为 \(2\)，短半轴平方为 \(1\)，所以焦距参数满足 \(c^2=2-1=1\)，即 \(F_1=(-1,0)\)，\(F_2=(1,0)\)，且 \(\lvert F_1F_2\rvert=2\).

若 \(PF_1=PF_2\)，则 \(P\) 在 \(y\) 轴上，椭圆与 \(y\) 轴交于 \((0,1)\)、\((0,-1)\)，有 \(2\) 个点.

若 \(PF_1=2\)，设 \(P=(x,y)\)，由椭圆方程得 \(y^2=1-\dfrac{x^2}{2}\)，于是 \((x+1)^2+y^2=4\)，即 \(x^2+4x-4=0\). 在 \([-\sqrt2,\sqrt2]\) 内只有 \(x=-2+2\sqrt2\)，对应 \(y=\pm\sqrt{1-\dfrac{x^2}{2}}\)，有 \(2\) 个点.

若 \(PF_2=2\)，则 \((x-1)^2+y^2=4\). 代入 \(y^2=1-\dfrac{x^2}{2}\) 得 \(x^2-4x-4=0\)，其根为 \(x=2\pm2\sqrt2\). 只有 \(x=2-2\sqrt2\in[-\sqrt2,\sqrt2]\) 合法，对应 \(y=\pm\sqrt{1-\dfrac{x^2}{2}}\)，有 \(2\) 个点. 三种情形互不重复，所以共有 \(2+2+2=6\) 个点.
\end{solution}

\begin{problem}
设 \(a_1\)，\(a_2\)，\(\cdots\)，\(a_6\) 为 \(1\)，\(2\)，\(3\)，\(4\)，\(5\)，\(6\) 的一个排列，则满足 \(\lvert a_1-a_2\rvert+\lvert a_3-a_4\rvert+\lvert a_5-a_6\rvert=3\) 的不同排列的个数为\fillinblank{}.
\end{problem}

\begin{answer}
\(48\)
\end{answer}

\begin{solution}
三个绝对值均为正整数且和为 \(3\)，所以必须分别等于 \(1\). 因而六个数只能配成 \((1,2)\)、\((3,4)\)、\((5,6)\) 这三个数对；这是因为差为 \(1\) 的三个不相交数对的唯一完全配对就是上述配对.

将这三个数对安排到 \((a_1,a_2)\)、\((a_3,a_4)\)、\((a_5,a_6)\) 的位置有 \(3!\) 种，每个数对内部交换次序有 \(2^3\) 种. 所以不同排列的个数为 \(3!\times2^3=48\).
\end{solution}

\begin{problem}
设 \(a\)，\(b\in\R\)，若函数 \(f(x)=x+\dfrac{a}{x}+b\) 在区间 \((1,2)\) 上有两个不同的零点，则 \(f(1)\) 的取值范围为\fillinblank{}.
\end{problem}

\begin{answer}
\((0,1)\)
\end{answer}

\begin{solution}
因为两个零点都在 \((1,2)\) 内，设它们为 \(r_1\)，\(r_2\)，且 \(1<r_1<r_2<2\). 由于 \(x\in(1,2)\) 时可以乘以 \(x\)，方程 \(f(x)=0\) 等价于 \(x^2+bx+a=0\)，所以 \(x^2+bx+a=(x-r_1)(x-r_2)\). 因此
\[
f(1)=1+a+b=(1-r_1)(1-r_2)=(r_1-1)(r_2-1)
\]
令 \(u=r_1-1\)，\(v=r_2-1\)，则 \(0<u<v<1\)，故 \(0<uv<1\). 反之，任取 \(0<t<1\)，取 \(u=\dfrac{2t}{1+t}\)、\(v=\dfrac{1+t}{2}\)，则 \(0<u<v<1\) 且 \(uv=t\). 通过 \(r_1=1+u\)、\(r_2=1+v\) 可构造出两个不同的零点，所以取值范围为 \((0,1)\).
\end{solution}

\section{单选题}

\begin{problem}
函数 \(f(x)=(x-1)^2\) 的单调递增区间是（\quad）.

\choices
    {\([0,+\infty)\)}
    {\([1,+\infty)\)}
    {\((-\infty,0]\)}
    {\((-\infty,1]\)}
\end{problem}

\begin{answer}
B
\end{answer}

\begin{solution}
函数图象为开口向上的抛物线，顶点为 \((1,0)\). 当 \(x\leq1\) 时，函数递减；当 \(x\geq1\) 时，函数递增. 因此单调递增区间是 \([1,+\infty)\)，故选 B.
\end{solution}

\begin{problem}
设 \(a\in\R\)，“\(a>0\)”是“\(\dfrac1a>0\)”的（\quad）条件.

\choices
    {充分非必要}
    {必要非充分}
    {充要}
    {既非充分也非必要}
\end{problem}

\begin{answer}
C
\end{answer}

\begin{solution}
若 \(a>0\)，则 \(a\neq0\)，且正数的倒数仍为正数，所以 \(\dfrac1a>0\). 反过来，若 \(\dfrac1a>0\)，则 \(a\neq0\)，并且 \(a\) 与其倒数同号，从而 \(a>0\). 两个条件可以互相推出，故为充要条件，选 C.
\end{solution}

\begin{problem}
过正方体中心的平面截正方体所得的截面中，不可能的图形是（\quad）.

\choices
    {三角形}
    {长方形}
    {对角线不相等的菱形}
    {六边形}
\end{problem}

\begin{answer}
A
\end{answer}

\begin{solution}
把正方体中心作为原点，正方体写成 \([-1,1]^3\). 过原点的截面关于原点中心对称，因此截面多边形的顶点成对出现，边数必须为偶数，不可能是三角形. 例如平面 \(z=0\) 截得正方形，正方形是长方形；平面 \(x+y+3z=0\) 的截面顶点为 \((1,1,-\frac23)\)、\((1,-1,0)\)、\((-1,-1,\frac23)\)、\((-1,1,0)\)，相邻边长均为 \(\dfrac{2\sqrt{10}}3\)，而两条对角线长分别为 \(\dfrac{2\sqrt{22}}3\)、\(2\sqrt2\)，所以它是对角线不相等的菱形；平面 \(x+y+z=0\) 截得六边形. 所以不可能的图形是三角形，选 A.
\end{solution}

\begin{problem}
如图，正八边形 \(A_1A_2\cdots A_8\) 的边长为 \(2\)，若 \(P\) 为该正八边形边上的动点，则 \(\overrightarrow{A_1A_3}\cdot \overrightarrow{A_1P}\) 的取值范围为（\quad）.

\bitmapfigure[float=inline,max width=0.28\linewidth]{img/2017/shanghai_spring/q16_fig1.png}

\choices
    {\([0,8+6\sqrt2]\)}
    {\([-2\sqrt2,8+6\sqrt2]\)}
    {\([-8-6\sqrt2,2\sqrt2]\)}
    {\([-8-6\sqrt2,8+6\sqrt2]\)}
\end{problem}

\begin{answer}
B
\end{answer}

\begin{solution}
取 \(A_1\) 为原点，令 \(A_1A_2\) 沿 \(x\) 轴正方向，边长为 \(2\). 依次取正八边形各边方向相差 \(45^\circ\)，则
\(A_1=(0,0)\)，\(A_2=(2,0)\)，\(A_3=(2+\sqrt2,\sqrt2)\)
且 \(\overrightarrow{A_1A_3}=(2+\sqrt2,\sqrt2)\). 各顶点坐标可依次写为
\(A_4=(2+\sqrt2,2+\sqrt2)\)，\(A_5=(2,2+2\sqrt2)\)，\(A_6=(0,2+2\sqrt2)\)
\(A_7=(-\sqrt2,2+\sqrt2)\)，\(A_8=(-\sqrt2,\sqrt2)\).
令 \(L(P)=\overrightarrow{A_1A_3}\cdot\overrightarrow{A_1P}\). 线性函数在正八边形边界上的最大、最小值出现在顶点. 逐点计算得
\(L(A_1)=L(A_7)=0\)，\(L(A_2)=L(A_6)=4+2\sqrt2\)
\(L(A_3)=L(A_5)=8+4\sqrt2\)，\(L(A_4)=8+6\sqrt2\)，\(L(A_8)=-2\sqrt2\).
所以取值范围为 \([-2\sqrt2,8+6\sqrt2]\)，故选 B.
\end{solution}

\section{解答题}

\begin{problem}
如图，长方体 \(ABCD-A_1B_1C_1D_1\) 中，\(AB=BC=2\)，\(AA_1=3\).

\begin{enumerate}
\item 求四棱锥 \(A_1-ABCD\) 的体积；
\item 求异面直线 \(A_1C\) 与 \(DD_1\) 所成角的大小.
\end{enumerate}

\bitmapfigure[max width=0.28\linewidth]{img/2017/shanghai_spring/q17_fig1.png}
\end{problem}

\begin{solution}
\begin{enumerate}
\item 底面 \(ABCD\) 是边长为 \(2\) 的正方形，面积为 \(4\)，且四棱锥的高为 \(AA_1=3\). 所以体积为 \(\dfrac13\times4\times3=4\).
\item 建立直角坐标系，取 \(A=(0,0,0)\)、\(B=(2,0,0)\)、\(C=(2,2,0)\)、\(D=(0,2,0)\)、\(A_1=(0,0,3)\)、\(D_1=(0,2,3)\). 直线 \(A_1C\) 与 \(DD_1\) 的方向向量可取 \(\bs u=(2,2,-3)\)、\(\bs v=(0,0,3)\). 令两异面直线所成的锐角为 \(\theta\)，则 \(\cos\theta=\dfrac{|\bs u\cdot\bs v|}{|\bs u||\bs v|}=\dfrac9{3\sqrt{17}}=\dfrac3{\sqrt{17}}\). 因而 \(\tan\theta=\dfrac{\sqrt{1-9/17}}{3/\sqrt{17}}=\dfrac{2\sqrt2}{3}\)，所以 \(\theta=\arctan\dfrac{2\sqrt2}{3}\).
\end{enumerate}
\end{solution}

\begin{problem}
设 \(a\in\R\)，函数 \(f(x)=\dfrac{2^x+a}{2^x+1}\).

\begin{enumerate}
\item 求 \(a\) 的值，使得 \(f(x)\) 为奇函数；
\item 若 \(f(x)<\dfrac{a+2}{2}\) 对任意 \(x\in\R\) 成立，求 \(a\) 的取值范围.
\end{enumerate}
\end{problem}

\begin{solution}
\begin{enumerate}
\item 令 \(t=2^x>0\)，则 \(f(x)=\dfrac{t+a}{t+1}\)，而 \(f(-x)=\dfrac{1+at}{t+1}\). 要使 \(f(-x)=-f(x)\) 对任意 \(t>0\) 成立，必须有 \(1+at=-t-a\)，即 \((a+1)(t+1)=0\). 因为 \(t+1>0\)，所以 \(a=-1\).
\item 仍令 \(t=2^x>0\). 不等式等价于
\[
\frac{a+2}{2}-\frac{t+a}{t+1}=\frac{at+2-a}{2(t+1)}>0
\]
因分母为正，只需 \(at+2-a>0\) 对任意 \(t>0\) 成立. 若 \(a<0\)，当 \(t\) 足够大时 \(at+2-a<0\)，不合题意；若 \(a=0\)，该式恒为 \(2>0\). 若 \(a>0\)，当 \(a\leq2\) 时 \(at+2-a>2-a\geq0\)，且因 \(t>0\) 实际为正；当 \(a>2\) 时，取足够小的正数 \(t<\dfrac{a-2}{a}\)，便有 \(at+2-a<0\). 因此 \(a\in[0,2]\).
\end{enumerate}
\end{solution}

\begin{problem}
某景区欲建造两条圆形观景步道 \(M_1\)，\(M_2\)（宽度忽略不计），如图，已知 \(AB\perp AC\)，\(AB=AC=AD=60\)（单位：米），要求圆 \(M_1\) 与 \(AB\)，\(AD\) 分别相切于点 \(B\)，\(D\)，圆 \(M_2\) 与 \(AC\)，\(AD\) 分别相切于点 \(C\)，\(D\).

\begin{enumerate}
\item 若 \(\angle BAD=60^\circ\)，求圆 \(M_1\)，\(M_2\) 的半径（结果精确到 \(0.1\) 米）；
\item 若步道 \(M_1\)，\(M_2\) 的每米造价分别为 \(0.8\text{ 千元}\) 和 \(0.9\text{ 千元}\)，如何设计圆 \(M_1\)，\(M_2\) 的大小，使总造价最低？最低总造价是多少？（结果精确到 \(0.1\text{ 千元}\)）.
\end{enumerate}

\bitmapfigure[max width=0.34\linewidth]{img/2017/shanghai_spring/q19_fig1.png}
\end{problem}

\begin{solution}
\begin{enumerate}
\item 设 \(\angle BAD=\theta\). 圆与两条相交直线相切时，圆心在角平分线上；连接圆心与切点，得到以半径为直角边、以 \(60\) 为邻直角边的直角三角形，所以 \(r=60\tan\dfrac{\theta}{2}\). 当 \(\theta=60^\circ\) 时，\(r_1=60\tan30^\circ=20\sqrt3\approx34.6\text{ 米}\). 又因为 \(AB\perp AC\)，有 \(\angle CAD=30^\circ\)，故 \(r_2=60\tan15^\circ=60(2-\sqrt3)\approx16.1\text{ 米}\).
\item 因为射线 \(AD\) 位于互相垂直的射线 \(AB\)、\(AC\) 之间，所以 \(0<\theta<90^\circ\). 设 \(u=\tan\dfrac{\theta}{2}\)，则 \(0<u<1\)，并且 \(r_1=60u\)，\(r_2=60\tan\left(45^\circ-\dfrac{\theta}{2}\right)=60\dfrac{1-u}{1+u}\). 两条圆形步道的总造价为
\[
C(u)=2\pi\left(0.8r_1+0.9r_2\right)=120\pi\left(0.8u+0.9\frac{1-u}{1+u}\right)
\]
令 \(g(u)=0.8u+0.9\dfrac{1-u}{1+u}\)，则 \(g'(u)=0.8-\dfrac{1.8}{(1+u)^2}\). 令 \(g'(u)=0\)，得 \((1+u)^2=\dfrac94\)，因 \(u>0\)，所以 \(u=\dfrac12\). 且 \(g'(u)<0\)（当 \(0<u<\dfrac12\)），\(g'(u)>0\)（当 \(\dfrac12<u<1\)），故此时总造价最低.

此时 \(r_1=60u=30\text{ 米}\)，\(r_2=60\dfrac{1-u}{1+u}=20\text{ 米}\)，最低总造价为 \(2\pi(0.8\times30+0.9\times20)=84\pi\approx263.9\text{ 千元}\).
\end{enumerate}
\end{solution}

\begin{problem}
已知双曲线 \(\Gamma:x^2-\dfrac{y^2}{b^2}=1\)（\(b>0\)），直线 \(l:y=kx+m\)（\(km\neq0\)），\(l\) 与 \(\Gamma\) 交于 \(P\)，\(Q\) 两点，\(P'\) 为 \(P\) 关于 \(y\) 轴的对称点，直线 \(P'Q\) 与 \(y\) 轴交于点 \(N(0,n)\).

\begin{enumerate}
\item 若点 \((2,0)\) 是 \(\Gamma\) 的一个焦点，求 \(\Gamma\) 的渐近线方程；
\item 若 \(b=1\)，点 \(P\) 的坐标为 \((-1,0)\)，且 \(\overrightarrow{NP'}=\dfrac32\,\overrightarrow{P'Q}\)，求 \(k\) 的值；
\item 若 \(m=2\)，求 \(n\) 关于 \(b\) 的表达式.
\end{enumerate}
\end{problem}

\begin{solution}
\begin{enumerate}
\item 双曲线的实半轴为 \(a=1\)，焦距参数满足 \(c^2=a^2+b^2=1+b^2\). 因为焦点 \((2,0)\) 在 \(x\) 轴上，所以 \(c=2\)，从而 \(b^2=3\). 双曲线的渐近线为 \(y=\pm\dfrac ba x\)，故方程为 \(y=\pm\sqrt3x\).
\item 当 \(b=1\) 且 \(P=(-1,0)\) 时，\(P'=(1,0)\). 设 \(Q=(x,y)\)，由向量条件 \(\overrightarrow{NP'}=\dfrac32\overrightarrow{P'Q}\)，得 \(N=P'-\dfrac32(Q-P')\). 因为 \(N\) 在 \(y\) 轴上，比较横坐标得 \(0=1-\dfrac32(x-1)\)，所以 \(x=\dfrac53\). 又 \(Q\) 在 \(x^2-y^2=1\) 上，故 \(y=\pm\dfrac43\). 直线 \(PQ\) 的斜率为 \(k=\dfrac{y-0}{\frac53-(-1)}=\pm\dfrac12\).
\item 设 \(P=(p,kp+2)\)、\(Q=(q,kq+2)\). 将直线 \(y=kx+2\) 代入双曲线，得 \((b^2-k^2)x^2-4kx-(b^2+4)=0\). 因此由韦达定理，\(p+q=\dfrac{4k}{b^2-k^2}\)，\(pq=-\dfrac{b^2+4}{b^2-k^2}\).

\(P'=(-p,kp+2)\)，所以直线 \(P'Q\) 在 \(y\) 轴上的截距为
\[
n=kp+2+\frac{k(q-p)}{p+q}p=2+\frac{2kpq}{p+q}=2-\frac{b^2+4}{2}=-\frac{b^2}{2}
\]
故 \(n=-\dfrac{b^2}{2}\).
\end{enumerate}
\end{solution}

\begin{problem}
已知函数 \(f(x)=\log_2\dfrac{1+x}{1-x}\).

\begin{enumerate}
\item 解方程 \(f(x)=1\)；
\item 设 \(x\in(-1,1)\)，\(a\in(1,+\infty)\)，证明：\(\dfrac{ax-1}{a-x}\in(-1,1)\)，且 \(f\!\left(\dfrac{ax-1}{a-x}\right)-f(x)=-f\!\left(\dfrac1a\right)\)；
\item 设数列 \(\{x_n\}\) 中，\(x_1\in(-1,1)\)，\(x_{n+1}=(-1)^{n+1}\dfrac{3x_n-1}{3-x_n}\)，\(n\in\N^*\)，求 \(x_1\) 的取值范围，使得 \(x_3\ge x_n\) 对任意 \(n\in\N^*\) 成立.
\end{enumerate}
\end{problem}

\begin{solution}
\begin{enumerate}
\item 函数有定义的条件为 \(-1<x<1\). 方程 \(f(x)=1\) 等价于 \(\dfrac{1+x}{1-x}=2\)，所以 \(1+x=2-2x\)，解得 \(x=\dfrac13\)，它满足定义域条件.
\item 记 \(y=\dfrac{ax-1}{a-x}\). 因为 \(a>1\)、\(-1<x<1\)，所以 \(a-x>0\). 又
\(1+y=\frac{(a-1)(1+x)}{a-x}>0\)，\(1-y=\frac{(a+1)(1-x)}{a-x}>0\)
故 \(-1<y<1\). 进一步有
\[
\frac{1+y}{1-y}=\frac{a-1}{a+1}\cdot\frac{1+x}{1-x}
\]
取以 \(2\) 为底的对数，得 \(f(y)=f(x)+\log_2\dfrac{a-1}{a+1}\). 而 \(f\left(\dfrac1a\right)=\log_2\dfrac{a+1}{a-1}\)，所以 \(f\left(\dfrac{ax-1}{a-x}\right)-f(x)=-f\left(\dfrac1a\right)\).
\item 记 \(T(x)=\dfrac{3x-1}{3-x}\)，由上一问知 \(T\) 将 \((-1,1)\) 映射到自身，且 \(f(T(x))=f(x)-f\left(\dfrac13\right)=f(x)-1\). 另外，直接由定义可得 \(f(-u)=-f(u)\)（\(-1<u<1\)）. 又
\(f'(x)=\frac{2}{(1-x^2)\ln 2}>0\)，\((-1<x<1)\)
所以 \(f\) 在 \((-1,1)\) 上严格递增. 令 \(y_n=f(x_n)\)，递推式化为
\[
y_{n+1}=\begin{cases}
y_n-1,&n\text{ 为奇数},\\
1-y_n,&n\text{ 为偶数}.
\end{cases}
\]
设 \(y_1=t\)，依次得到 \(y_2=t-1\)，\(y_3=2-t\)，\(y_4=1-t\)，\(y_5=t\)，故之后以四项为周期. 由于 \(f\) 严格递增，题设条件等价于 \(y_3\ge y_n\) 对任意 \(n\) 成立. 在四个值中，\(y_3\ge y_4\) 恒成立，\(y_3\ge y_2\) 等价于 \(t\le\dfrac32\)，而 \(y_3\ge y_1\) 等价于 \(t\le1\)，后者更严格，因此只需 \(t\le1\). 又 \(f\left(\dfrac13\right)=1\)，故 \(t\le1\) 等价于 \(x_1\le\dfrac13\). 联合 \(x_1\in(-1,1)\)，得到 \(x_1\in\left(-1,\dfrac13\right]\).
\end{enumerate}
\end{solution}
